\subsection{单射或满射线性映射的扰动}

定理 1. --- 设 E 和 F 是分离局部凸空间，u 是从 E 到 F 中的严格态射，其核有限维且其像是闭的。设 h 是从 E 到 F 中的紧线性映射。线性映射 \(u + h\) 是严格态射，它的核有限维，且它的像是闭的。

由 III, p. 71 的命题 1，存在 E 中 0 的一个闭邻域 V，使得 \(u\) 在 V 上的限制是真映射。由于 \(h\) 是紧的，存在包含于 V 的 0 的闭邻域 W，使得集 \(h(\mathrm{W})\) 是相对紧的。令 \(\mathrm{C} = \overline{h(\mathrm{W})}\) 。\(u\) 在 W 上的限制是真映射（TG, I, p. 74, 推论 1）。映射 \(\alpha: x \mapsto (u(x), h(x))\) （从 W 到 \(\mathrm{F} \times \mathrm{C}\) 中）是真映射，因为映射 \(\mathrm{pr}_1 \circ \alpha = u|\mathrm{W}\) （从 W 到 F 中）是真映射（TG, I, p. 73, 命题 5, \(d\) ）。映射 \(\beta: (y, z) \mapsto (y + z, z)\) （从 \(\mathrm{F} \times \mathrm{C}\) 到 \(\mathrm{F} \times \mathrm{C}\) 中）是同胚，而映射 \(\mathrm{pr}_1\) （从 \(\mathrm{F} \times \mathrm{C}\) 到 F 中）是真映射（TG, I, p. 77, 推论 5）。于是复合映射 \(\mathrm{pr}_1 \circ \beta \circ \alpha\) （从 W 到 F 中）是真映射（TG, I, p. 73, 命题 5, \(a\) ）。但这个映射正是 \(u + h\) 在 W 上的限制。由 III, p. 71 的命题 1，\(u + h\) 是严格态射，它的核有限维，且它的像是闭的。

引理 2. --- 设 E 和 F 是弗雷歇空间，且 \(u \in \mathcal{L}(\mathrm{E};\mathrm{F})\) 。下列条件等价：

\begin{enumerate}
\def\labelenumi{(\roman{enumi})}
\item
  u 的余核有限维；
\item
  \(^{t}u \colon \mathrm{F}_c' \to \mathrm{E}_c'\) 是具有闭像且核有限维的严格态射。
\item
  \(\Longrightarrow\) (ii)：假设 u 的余核有限维。此时映射 u 是严格态射（III, p. 52 的引理 6）。由 EVT, IV, p. 27 的命题 2，\(^{t}u\) 的像在赋有弱拓扑的 E' 中是闭的，从而在 \(E_{c}'\) 中更是闭的。\(^{t}u\) 的核是 u 的像的正交补（同前）；故它有限维。最后，由 EVT, IV, p. 28 的定理 1，\(^{t}u\) 是从 \(F_{c}'\) 到 \(E_{c}'\) 中的严格态射。
\item
  \(\Longrightarrow\) (i)：假设 \(^t u\colon F_c' \to E_c'\) 是严格态射。由 EVT, IV, p. 28 的定理 1，\(u\) 的像是闭的。\(^t u\) 的核是 \(\operatorname{Im}(u)\) 的正交补（依照 EVT, IV, p. 27, 命题 2）；若这个核有限维，则 \(u\) 的像在 F 中有有限余维数。
\end{enumerate}

定理 2. --- 设 E 和 F 是弗雷歇空间，u: E → F 是余核有限维的连续线性映射，h: E → F 是紧线性映射。线性映射 u + h 是严格态射，它的余核有限维，且它的像是闭的。

由 III, p. 52 的引理 6，只需证明 \(u + h\) 的余核有限维。而由 III, p. 9 的命题 9，映射 \(^{t}h\) （从 \(F_{c}^{\prime}\) 到 \(E_{c}^{\prime}\) 中）是紧的。定理 2 由定理 1 与引理 2 得出。

